\section*{NeurIPS Paper Checklist}

\begin{enumerate}

\item {\bf Claims}
    \item[] Question: Do the main claims made in the abstract and introduction accurately reflect the paper's contributions and scope?
    \item[] Answer: \answerYes{}
    \item[] Justification: The abstract and introduction enumerate exactly the contributions delivered: a five-encoder stack (\S\ref{sec:method}), Perceiver IO fusion (\S\ref{sec:fusion}), a PPO cascade controller (\S\ref{sec:cascade}), the SENTINEL-DB database (\S\ref{sec:db}), and a 20-experiment evaluation suite (\S\ref{sec:exp}). All quantitative headline claims (full-fusion AUROC~$0.9919$, $95\%$ CI $[0.9907,0.9934]$, paired permutation $p\!=\!1.6{\times}10^{-3}$, FPR $0.000$ across $10$ NEON reference sites, $12.8\times$ signal-to-noise) are reported with the source experiment file in the released repository.

\item {\bf Limitations}
    \item[] Question: Does the paper discuss the limitations of the work performed by the authors?
    \item[] Answer: \answerYes{}
    \item[] Justification: \S\ref{sec:limits} discusses (i)~the behavioral channel's perfect AUROC reflecting controlled EPA ECOTOX assays rather than field deployment, (ii)~advisory-date label noise in the historical benchmark, (iii)~limits of multispectral inversion at $10$\,m for TSS/nitrate, (iv)~ortholog-coverage constraints on cross-species ToxiGene transfer, and (v)~North-American/European bias in SENTINEL-DB. Appendix~\ref{app:conformal} additionally reports honestly that conformal coverage is met only on the high-volume behavioral channel.

\item {\bf Theory assumptions and proofs}
    \item[] Answer: \answerNA{}
    \item[] Justification: The paper is empirical. The only formal guarantees we appeal to (distribution-free conformal coverage) are inherited unmodified from \citep{vovk2005conformal,angelopoulos2021conformal}; we do not state or prove new theorems.

\item {\bf Experimental result reproducibility}
    \item[] Question: Does the paper fully disclose all the information needed to reproduce the main experimental results of the paper to the extent that it affects the main claims and/or conclusions of the paper?
    \item[] Answer: \answerYes{}
    \item[] Justification: All training data are derived from public sources (USGS NWIS, EPA WQP/ECOTOX, NEON, ESA Sentinel-2, EMP, NCBI GEO, GBIF, GRQA) referenced in \S\ref{sec:db}. The released codebase contains data-acquisition scripts, training scripts for each encoder and the fusion/controller, configuration files, and the JSON outputs of every experiment in \texttt{results/}. Hyperparameters are summarised in Appendix~\ref{app:hp}.

\item {\bf Open access to data and code}
    \item[] Answer: \answerYes{}
    \item[] Justification: Code and SENTINEL-DB harmonisation pipeline are released under MIT. An anonymised snapshot is included with the submission's supplementary material; raw upstream data are pulled from the public sources above using the provided scripts.

\item {\bf Experimental setting/details}
    \item[] Answer: \answerYes{}
    \item[] Justification: \S\ref{sec:method} specifies architectures and pre-training/fine-tuning recipes per encoder. Optimizer, learning-rate, batch-size, and schedule details are in Appendix~\ref{app:hp} and \texttt{configs/default.yaml}. Train/val/test splits and per-experiment evaluation protocols are given in \S\ref{sec:exp} and Appendix~\ref{app:baselines}.

\item {\bf Experiment statistical significance}
    \item[] Answer: \answerYes{}
    \item[] Justification: Headline gain over the best single modality uses a paired permutation test ($n_\text{perm}\!=\!10{,}000$, $p\!=\!1.6{\times}10^{-3}$). Per-modality and fusion AUROCs are reported with $95\%$ bootstrap confidence intervals. Modality ablation contributions are averaged over $15$ pairwise comparisons per modality. Robustness to missing modalities uses $100$ random-drop trials; calibration uses $50$ MC-dropout passes. The factors of variability captured (window draws within events; modality-mask draws; MC-dropout sampling) are stated alongside each number.

\item {\bf Experiments compute resources}
    \item[] Answer: \answerYes{}
    \item[] Justification: Each encoder pre-training and fine-tuning run fits on a single NVIDIA A100 (40\,GB) within $<\!36$\,h; fusion training requires a single A100 for $\sim\!12$\,h; PPO cascade training runs $500$k timesteps with $8$ CPU rollout workers feeding a single A100, completing in $\sim\!12$\,h. Total project compute, including failed experiments and benchmark sweeps not appearing in the paper, is approximately $4$ A100-weeks.

\item {\bf Code of ethics}
    \item[] Answer: \answerYes{}
    \item[] Justification: The work uses only public, non-personal environmental and ecotoxicology data; no human subjects are involved; the released artefacts, including SENTINEL-DB, are explicitly intended for water-safety research and respect upstream licenses.

\item {\bf Broader impacts}
    \item[] Answer: \answerYes{}
    \item[] Justification: \S\ref{sec:impact} discusses (i)~positive impacts on public health and environmental monitoring, (ii)~dual-use risk of evading detection (mitigated by orthogonal modality cascade), (iii)~over-reliance risk (mitigated by published conformal coverage and explicit recommendation against unilateral regulatory use), and (iv)~equity risk in sensor placement (mitigated by an optional vulnerability-weighting hook in the released code).

\item {\bf Safeguards}
    \item[] Answer: \answerYes{}
    \item[] Justification: The released models are detection-oriented and present low direct misuse risk. We nevertheless (i)~require all five-modality cascade activation for high-stakes alerts, (ii)~publish conformal coverage and calibration diagnostics so that downstream users can quantify reliability, (iii)~document the dual-use considerations in the model card, and (iv)~discourage unilateral regulatory use without confirmatory laboratory analysis.

\item {\bf Licenses for existing assets}
    \item[] Answer: \answerYes{}
    \item[] Justification: All upstream datasets are cited (USGS NWIS \citep{usgs_nwis}, EPA WQP \citep{epa_wqp}, EPA ECOTOX \citep{epa_ecotox}, NEON \citep{neon}, GRQA v1.3 \citep{virro2021grqa}, EMP 16S \citep{thompson2017earthmicrobiome}, ESA Sentinel-2 \citep{drusch2012sentinel2}, NCBI GEO, GBIF). All are public-domain or released under permissive licenses (US federal works are public domain; ESA Sentinel data are released under the Copernicus Open Access policy; EMP and GEO are public). Original licenses are reproduced in the supplementary \texttt{LICENSES.md}.

\item {\bf New assets}
    \item[] Answer: \answerYes{}
    \item[] Justification: SENTINEL-DB and the trained model checkpoints are released under MIT. The supplementary material includes a data sheet for SENTINEL-DB (sources, harmonisation steps, quality tiers, known biases) and a model card per encoder (training data, intended use, known failure modes, evaluation results). Anonymised at submission time.

\item {\bf Crowdsourcing and research with human subjects}
    \item[] Answer: \answerNA{}
    \item[] Justification: The work involves no crowdsourcing and no research with human subjects.

\item {\bf Institutional review board (IRB) approvals}
    \item[] Answer: \answerNA{}
    \item[] Justification: The work involves no human subjects research.

\item {\bf Declaration of LLM usage}
    \item[] Answer: \answerNA{}
    \item[] Justification: Large language models were used only for incidental writing and code-style polish; they are not a component of the core methodology, encoders, fusion layer, controller, evaluation, or any reported result.

\end{enumerate}
